中. 然而,在电压和气流的 $99\%$ 限制的椭圆质量控制图(\hyperref[fig:5.10]{图 5.10})中,可以看到 31 号失控,原因是气流过大. 在 $\ln(\text{气流})$ 的单变量的 $\bar{X}$ 控制图(\hyperref[fig:5.11]{图 5.11})中,我们可以看到这个点在 $3\sigma$ 控制限之外. 看起来 32 号的气流数据被重置了. 在所有其他的单变量 $\bar{X}$ 控制图中,所有的点都在 $3\sigma$ 控制限之内.

\begin{mvfigure}
\begin{minipage}{0.5\linewidth}\centering
\begin{tikzpicture}[x=0.155cm,y=0.16cm]
  \draw[->] (0,0) -- (42,0) node[right,font=\scriptsize]{序号};
  \draw[->] (0,0) -- (0,15.5) node[above,font=\scriptsize]{$T^{2}$};
  \foreach \y in {0,2,4,6,8,10,12,14} \draw (0,\y) -- (-1,\y) node[left,font=\tiny]{\y};
  \draw[dashed] (0,9.5) -- (41,9.5);
  \node[font=\tiny,anchor=south east] at (41,9.5) {$95\%$限制};
  \draw (0,12) -- (41,12);
  \node[font=\tiny,anchor=south east] at (41,12) {$99\%$限制};
  \draw (1,2)--(2,3)--(3,5)--(4,6)--(5,4)--(6,3)--(7,7)--(8,6)--(9,4)--(10,3)--(11,5)--(12,2)--(13,8)--(14,6)--(15,4)--(16,7)--(17,3)--(18,2)--(19,5)--(20,4)--(21,3)--(22,7)--(23,6)--(24,4)--(25,3)--(26,5)--(27,2)--(28,4)--(29,3)--(30,2)--(31,10.5)--(32,3)--(33,2)--(34,4)--(35,3)--(36,2)--(37,3)--(38,2)--(39,1)--(40,2);
  \foreach \x/\y in {1/2,2/3,3/5,4/6,5/4,6/3,7/7,8/6,9/4,10/3,11/5,12/2,13/8,14/6,15/4,16/7,17/3,18/2,19/5,20/4,21/3,22/7,23/6,24/4,25/3,26/5,27/2,28/4,29/3,30/2,31/10.5,32/3,33/2,34/4,35/3,36/2,37/3,38/2,39/1,40/2}
    \fill (\x,\y) circle (1.1pt);
\end{tikzpicture}
\phantomsection\label{fig:5.9}\mvcap{图 5.9\quad 焊接数据的 $95\%$ 和 $99\%$ 的 $T^{2}$ 控制图}
\end{minipage}\hfill
\begin{minipage}{0.46\linewidth}\centering
\begin{tikzpicture}[x=1.15cm,y=1.4cm]
  \draw[->] (0,0) -- (3.9,0) node[right,font=\scriptsize]{电压};
  \draw[->] (0,0) -- (0,2.3) node[above,font=\scriptsize]{$\ln$(气流)};
  \foreach \x/\lab in {0/20.5,0.5/21.0,1.0/21.5,1.5/22.0,2.0/22.5,2.5/23.0,3.0/23.5,3.5/24.0}
    {\draw (\x,0) -- (\x,-0.06); \node[below,font=\tiny] at (\x,-0.08) {\lab};}
  \foreach \y/\lab in {0/3.85,0.5/3.90,1.0/3.95,1.5/4.00,2.0/4.05}
    {\draw (0,\y) -- (-0.06,\y); \node[left,font=\tiny] at (-0.06,\y) {\lab};}
  \draw (1.8,1.1) ellipse (1.05 and 0.55);
  \foreach \x/\y in {1.5/1.0,1.6/1.2,1.7/0.9,1.8/1.3,1.9/1.0,2.0/1.2,2.1/0.9,2.2/1.1,1.6/0.8,1.8/1.4,2.0/0.7,2.2/1.3,1.4/1.1,1.9/0.8,2.1/1.4,1.7/1.3,2.3/1.1,1.5/0.9,1.8/0.7,2.0/1.5,1.3/1.0,1.9/1.2,2.1/0.8,1.6/1.1,1.8/1.2,2.0/1.1,2.2/0.9,1.7/0.8,2.4/1.0,1.9/1.0}
    \fill (\x,\y) circle (0.9pt);
  \fill (1.7,1.9) circle (0.9pt);
\end{tikzpicture}
\phantomsection\label{fig:5.10}\mvcap{图 5.10\quad $\ln$(气流)和电压的 $99\%$ 椭圆质量控制图}
\end{minipage}
\end{mvfigure}

在这个例子中,通过组合到一个 $T^{2}$ 值里,单个变量的变化被掩盖在 $99\%$ 的限制里,或者几乎在 $95\%$ 的限制里被掩盖了.

\begin{mvfigure}
\begin{tikzpicture}[x=0.155cm,y=1.5cm]
  \draw[->] (0,0) -- (42,0) node[right,font=\scriptsize]{序号};
  \draw[->] (0,0) -- (0,2.7) node[above,font=\scriptsize]{$\ln$(气流)};
  \foreach \y/\lab in {0.4/3.90,1.4/3.95,2.4/4.00} \draw (0,\y) -- (-1,\y) node[left,font=\tiny]{\lab};
  \draw (0,2.5) -- (41,2.5) node[right,font=\tiny]{UCL=4.005};
  \draw (0,0.32) -- (41,0.32) node[right,font=\tiny]{LCL=3.896};
  \draw (0,1.42) -- (41,1.42) node[right,font=\tiny]{均值=3.951};
  \draw (1,1.3)--(2,1.5)--(3,1.2)--(4,1.6)--(5,1.1)--(6,1.4)--(7,1.7)--(8,1.2)--(9,1.5)--(10,1.3)--(11,1.6)--(12,1.2)--(13,1.4)--(14,1.1)--(15,1.5)--(16,1.3)--(17,1.6)--(18,1.2)--(19,1.4)--(20,1.5)--(21,1.2)--(22,1.6)--(23,1.3)--(24,1.5)--(25,1.2)--(26,1.4)--(27,1.6)--(28,1.3)--(29,1.5)--(30,1.2)--(31,2.5)--(32,1.4)--(33,1.3)--(34,1.5)--(35,1.2)--(36,1.4)--(37,1.3)--(38,1.5)--(39,1.2)--(40,1.4);
  \foreach \x/\y in {1/1.3,2/1.5,3/1.2,4/1.6,5/1.1,6/1.4,7/1.7,8/1.2,9/1.5,10/1.3,11/1.6,12/1.2,13/1.4,14/1.1,15/1.5,16/1.3,17/1.6,18/1.2,19/1.4,20/1.5,21/1.2,22/1.6,23/1.3,24/1.5,25/1.2,26/1.4,27/1.6,28/1.3,29/1.5,30/1.2,31/2.5,32/1.4,33/1.3,34/1.5,35/1.2,36/1.4,37/1.3,38/1.5,39/1.2,40/1.4}
    \fill (\x,\y) circle (1.1pt);
\end{tikzpicture}
\phantomsection\label{fig:5.11}\mvcap{图 5.11\quad $\ln$(气流)的单变量的 $\bar{X}$ 控制图}
\end{mvfigure}

\mvsub{未来单个观察值的控制域}

现在的目标是:当一个过程稳定时,利用收集到的数据 $\mathbf{x}_1,\mathbf{x}_2,\cdots,\mathbf{x}_n$ 来确定某一个或某些未来观测值的控制域. 预期未来观测值将落入其内的区域称为预测域. 若过程稳定,则观测值可视为独立同分布的,其分布服从 $N_p(\boldsymbol{\mu},\boldsymbol{\Sigma})$. 由于这些控制域除了监控质量之外还具有更普遍的重要性,我们在结论 5.6 中先给出其基本的分布理论.

\mvthm{结论 5.6}{设 $\mathbf{X}_1,\mathbf{X}_2,\cdots,\mathbf{X}_n$ 独立同分布,分布为 $N_p(\boldsymbol{\mu},\boldsymbol{\Sigma})$,记 $\mathbf{X}$ 为来自相同分布的一个未来观测值,此时
\[
T^{2}=\frac{n}{n+1}(\mathbf{X}-\bar{\mathbf{X}})'\mathbf{S}^{-1}(\mathbf{X}-\bar{\mathbf{X}})\ \text{服从分布}\ \frac{(n-1)p}{n-p}F_{p,n-p}
\]
且其 $100(1-\alpha)\%$ 的 $p$ 维预测椭球由满足如下不等式的全部 $\mathbf{x}$ 决定:}
